% TOP_PROBLEM: {"id":30007581,"problem_number":"LOCAL-30007581","title":"Adaptive market maintenance","category_id":21,"category":{"id":21,"name":"economics","display_name":"Economics","description":"Open economic theory statements, published research agendas, and proposed scoped specifications; question provenance and historical priority are qualified per record.","slug":"economics","order_index":21,"created_at":"2026-10-08T00:00:00Z"},"status":"provenance_pending","external_url":null,"legacy_ids":[],"published":false,"statement":"Determine attainable regret for adaptive market rules when exogenous conditions vary and each candidate rule sequence induces its own participation and adaptation trajectory, under the specified welfare, feedback, and switching constraints.","economics":{"original_id":70,"field":"Mechanism design and market design","kind":"design","formalization_basis":"proposed_specification","source_status":"not_verified","priority_status":"not_established","priority_note":"No primary statement of the exact online regret and endogenous-participation problem was verified; the working paper is supporting motivation.","earliest_verified_reference":null,"current_reference":{"authors":["Alvin E. Roth"],"title":"Market Design and Maintenance","year":2023,"url":"https://www.nber.org/papers/w31947","doi":"10.3386/w31947","locator":"Abstract (working paper31947); introduction in July17,2024 chapter revision","evidence_summary":"The paper explicitly treats adaptation to a changing larger market as an ongoing design requirement; it does not state our online-regret formulation."},"formalization_note":"Proposed online-maintenance formulation supported only by a broad maintenance framework. Exogenous conditions are shared across policies, while each policy induces its own counterfactual participant-state trajectory. It remains unverified as a canonical open problem."},"statusQualification":"Provenance pending: an explicit published statement of this exact open question was not verified. This is a catalogue-authored candidate specification, not a certified open conjecture. Proposed online-maintenance formulation supported only by a broad maintenance framework. Exogenous conditions are shared across policies, while each policy induces its own counterfactual participant-state trajectory. It remains unverified as a canonical open problem.","statusReviewedAt":"2026-10-08"}
% FIRST_POSTED: 2026-10-08
% ENTRY_KIND: statement_only
% !TeX program = lualatex
\documentclass[11pt]{article}
\usepackage{fontspec}
\usepackage[a4paper,margin=25mm]{geometry}
\usepackage{amsmath,amssymb,mathtools}
\usepackage{xurl}
\usepackage[hidelinks]{hyperref}
\newcommand{\sref}[2]{\hyperlink{source:#1:#2}{[S#2]}}
\setlength{\emergencystretch}{3em}
\begin{document}
% problemId:30007581
\section*{Adaptive market maintenance}\label{30007581}
\subsection{Definitions and mathematical statement}
A market operator repeatedly chooses a rule \(m_t\) from a finite set \(\mathcal M\), covering application limits, timing, signaling, or allocation priorities. Let \(\theta_t\) be exogenous market conditions; their evolution does not depend on selected rules. A participant state \(s_t\), including participation and adaptation, evolves as \(s_{t+1}=f(s_t,m_t,\theta_t)\), for a specified transition function \(f\) and common initial state \(s_1\). Strategic responses are incorporated into \(f\) and welfare \(W(m_t,s_t,\theta_t)\in[0,1]\). The operator observes outcomes under its chosen rule and incurs switching cost \(c(m_{t-1},m_t)\geq0\). For every comparator rule sequence \(m^*\), define its own counterfactual states by \(s^*_{t+1}=f(s^*_t,m^*_t,\theta_t)\), with \(s^*_1=s_1\). Set
\[
 R_T=\max_{m^*\text{ feasible}}\sum_t[W(m^*_t,s^*_t,\theta_t)-c(m^*_{t-1},m^*_t)]
 -\mathbb E\sum_t[W(m_t,s_t,\theta_t)-c(m_{t-1},m_t)].
\]
Both sequences obey the same announced switching budget and initial rule. For exogenous sequences satisfying \(\sum_{t=2}^T d(\theta_t,\theta_{t-1})\leq V\), with specified metric \(d\), determine minimax attainable regret over feasible adaptive policies, under explicit assumptions on state persistence, feedback, and strategic-response identification. Identify conditions allowing sublinear regret and failures caused by irreversible participation changes or unobserved counterfactuals. Specify the information needed to estimate welfare and transitions rather than treating observed outcomes as unbiased evaluations of alternatives. The maintenance source motivates adaptation as circumstances change; this online control specification was not verified as a published open problem.

\par\textbf{Formulation and scope.} Proposed online-maintenance formulation supported only by a broad maintenance framework. Exogenous conditions are shared across policies, while each policy induces its own counterfactual participant-state trajectory. It remains unverified as a canonical open problem.

\par\textbf{Open-question provenance.} An explicit published open-question statement was not verified. Sources below are supporting literature and must not be treated as a first listing.

\par\textbf{Historical priority.} No primary statement of the exact online regret and endogenous-participation problem was verified; the working paper is supporting motivation.

\subsection{Short English statement}
Determine attainable regret for adaptive market rules when exogenous conditions vary and each candidate rule sequence induces its own participation and adaptation trajectory, under the specified welfare, feedback, and switching constraints.

\subsection{Sources}
\begin{itemize}\raggedright
\item[\textnormal{[S1]}]\hypertarget{source:30007581:1}{} Alvin E. Roth. 2023. Market Design and Maintenance. Abstract (working paper31947); introduction in July17,2024 chapter revision.
\url{https://www.nber.org/papers/w31947}.
DOI: 10.3386/w31947.
{\small\textit{Supporting or current literature; see evidence qualification. The paper explicitly treats adaptation to a changing larger market as an ongoing design requirement; it does not state our online-regret formulation.}}
\end{itemize}
\end{document}
